\scanpage{14}
\textbf{C. R. Acad. Sc. Paris, t. 278 (24 juin 1974)}\hfill\textbf{Série A — 1599}

\rule{\linewidth}{0.4pt}

ALGÈBRE. — \emph{Formes quadratiques, algèbres de Clifford et signatures.}

Note (*) de M. \textbf{Olivier Laborde}, présentée par M. René Garnier.

Détermination du noyau de l'homomorphisme du groupe de Witt d'un corps de nombres ou
d'un anneau d'entiers algébriques dans le groupe de Brauer-Wall. Le calcul du groupe de Witt
est ramené à celui de $\mathrm K^*/\mathrm K^{*2}$ pour un corps de nombres et à celui du groupe des extensions
quadratiques pour un anneau d'entiers algébriques.

1. \textsc{Introduction}. — Étant donné un anneau commutatif unitaire A, il existe un homomorphisme
de groupe $f_{\mathrm A}$ de $\operatorname{Witt}(\mathrm A)$ dans $\operatorname{BW}(\mathrm A)$ définit dans (${}^2$) comme suit : si $q$ est
une forme quadratique non dégénérée sur un A-module projectif de type fini P, on note $[q]$
sa classe d'équivalence dans $\operatorname{Witt}(\mathrm A)$ et si D est une A-algèbre d'Azumaya graduée,
on note $[\mathrm D]$ sa classe d'équivalence dans $\operatorname{BW}(\mathrm A)$; alors $f_{\mathrm A}([q])=[\mathrm C(q)]$. L'image
de cet homomorphisme s'appelle le groupe de Hasse-Wall de A et on le note $\operatorname{HW}(\mathrm A)$;
le caractère fonctoriel de Witt, BW, HW et $f$ sera utilisé. D'après (${}^6$), si A est un anneau
de Dedekind dont le corps de fractions est tel que toute forme quadratiques de plus de
cinq variables soit isotrope, les groupes $\operatorname{Witt}(\mathrm A)$ et $\operatorname{HW}(\mathrm A)$ sont isomorphes.

En outre, dans (${}^9$), il est démontré que si K est un corps tel que les groupes $\operatorname{Witt}(\mathrm K)$
et $\operatorname{HW}(\mathrm K)$ soient isomorphes, alors le corps K possède la propriété indiquée ci-dessus.

Dans ce qui suit, nous déterminerons le noyau $\operatorname{LW}(\mathrm A)$ de l'homomorphisme $f_{\mathrm A}$ lorsque A
est un corps de nombres algébriques ou l'anneau de ses entiers et nous indiquerons à quels
invariants se ramène le calcul du groupe de Hasse-Wall.

Ce travail a été inspiré par (${}^7$).

2. \textsc{Les corps de nombres}. — Soit K un corps de nombre, i. e. une extension algébrique
finie du corps Q des nombres rationnels. Nous partagerons l'ensemble $\Omega$ des valuations
de K de la façon suivante :

(i) \emph{Les valuations réelles} (en nombre fini) : Ces valuations $(r_i)_{1\le i\le r}$ définissent sur K
une structure d'espace topologique et les complétés $\mathrm K_{r_i}$ sont isomorphes au corps des
nombres réels.

(ii) \emph{Les valuations complexes et les valuations discrètes} (en nombre infini) : Ces
valuations $(c_j)_{j\in\mathrm J}$ définissant sur K une structure d'espace topologique dont les complétés $\mathrm K_{c_j}$
sont isomorphes dans le cas des valuations complexes au corps des nombres complexes
et dans le cas des valuations discrètes à un corps local de caractéristique nulle, c'est-à-dire
dans ce cas une extension finie d'un corps $p$-adique.

A chaque valuation $v$ de K correspond une application de K dans $\mathrm K_v$. Le caractère
fonctoriel des groupes Witt et HW nous conduit au diagramme

\diagram{\begin{tikzcd}[ampersand replacement=\&,column sep=small]
0\arrow[r]\&\operatorname{LW}(\mathrm K)\arrow[r]\arrow[d,"\alpha"']\&
\operatorname{Witt}(\mathrm K)\arrow[r]\arrow[d,"\beta"']\&
\operatorname{HW}(\mathrm K)\arrow[r]\arrow[d,"\gamma"']\&0\\
0\arrow[r]\&\displaystyle\prod_{v\in\Omega}\operatorname{LW}(\mathrm K_v)\arrow[r]\&
\displaystyle\prod_{v\in\Omega}\operatorname{Witt}(\mathrm K_v)\arrow[r]\&
\displaystyle\prod_{v\in\Omega}\operatorname{HW}(\mathrm K_v)\arrow[r]\&0
\end{tikzcd}}

\textbf{C. R., 1974, 1\textsuperscript{er} \emph{Semestre}. (T. 278, N\textsuperscript{o} 26)}\hfill\textbf{Série A — 115}

\mspage{13}
